
This study tested whether linguistic agency markers validated in human psychology could serve as computational proxies for measuring agency preservation in RLHF evaluation. Analyzing 169,352 preference pairs from the HH-RLHF dataset, we found that while markers are reliably extractable (precision=100\%, CV=0.781), they fail comprehensive construct validation: effect sizes are 88\% below meaningful thresholds (d=-0.018 vs. required $d\geq 0.15$ despite p<0.001), internal consistency is 40\% below acceptable levels ($\alpha =0.42$ vs. required $\alpha$>0.7), and cross-split replication fails (0/2 splits passed).

The findings establish three contributions: (1) proxy validity requires multi-criterion empirical validation rather than theoretical assumption based on face validity, (2) statistical significance (p<0.001) achieved through large samples (N=169K) does not guarantee practical relevance without effect size thresholds, and (3) measurement feasibility (reliable extraction) is distinct from construct validity (measuring the intended construct).

The negative result prevents premature deployment of invalid automated metrics while establishing validation requirements for future computational operationalization of bidirectional alignment. Until validated proxies exist, agency preservation monitoring requires direct user studies or alternative measurement approaches with demonstrated construct validity.
